\section{Continuous Assessment 1 (CA-1) --- Set B}
\label{sec:ca1_set_b}

\LPUPaperHeader{Continuous Assessment (CA-1)}{B}{50 Minutes}{30 Marks}{CO1 (Matrices \& Systems of Linear Equations), CO2 (Higher-Order ODEs \& Calculus)}

\begin{center}
\sffamily\textbf{\color{AlertRed}Note: Attempt all six questions. Each question carries 5 marks.}
\end{center}

% ------------------------------------------------------------------------------
% QUESTION PAPER
% ------------------------------------------------------------------------------

\questionitem{1}{
    Find for what values of $\lambda$ and $\mu$ the system of linear equations:
    \[
    \begin{aligned}
    x + y + z &= 6 \\
    x + 2y + 3z &= 10 \\
    x + 2y + \lambda z &= \mu
    \end{aligned}
    \]
    has: \textbf{(i)} a unique solution, \textbf{(ii)} no solution, and \textbf{(iii)} infinitely many solutions.
}{5}{CO1}{L3 Apply}

\questionitem{2}{
    Solve the third-order linear homogeneous differential equation by the operator method:
    \[
    y''' + y = 0 \quad \iff \quad (D^3 + 1)y = 0
    \]
}{5}{CO2}{L3 Apply}

\questionitem{3}{
    Find the rank of the $4\times 4$ matrix $A$, given that parameters $w, r, s, t \neq 0$:
    \[
    A = \begin{bmatrix}
    1 & w & r & 0 \\
    0 & s & t & 1 \\
    -1 & -w & -r & 0 \\
    0 & s & t & 1
    \end{bmatrix}
    \]
}{5}{CO1}{L2 Understand}

\questionitem{4}{
    Verify the \textbf{Cayley-Hamilton Theorem} for the matrix:
    \[
    A = \begin{bmatrix} 1 & 2 \\ 3 & 4 \end{bmatrix}
    \]
    and hence compute $A^{-1}$ and $A^3$.
}{5}{CO1}{L3 Apply}

\questionitem{5}{
    Determine whether the matrix:
    \[
    A = \begin{bmatrix} 2 & 3+4i \\ 3-4i & 5 \end{bmatrix}
    \]
    is a \textbf{Hermitian matrix}. Find its eigenvalues and verify that they are strictly real.
}{5}{CO1}{L3 Apply}

\questionitem{6}{
    Find the complete general solution of the second-order differential equation:
    \[
    (D^2 - 4D + 4)y = e^{2x}
    \]
}{5}{CO2}{L3 Apply}

\vspace{5mm}
\hrule
\vspace{5mm}

% ------------------------------------------------------------------------------
% COMPLETE STEP-BY-STEP SOLUTIONS
% ------------------------------------------------------------------------------
\subsection*{Solutions \& Marking Scheme (CA-1 Set B)}

% Solution 1
\begin{solutionbox}[Solution to Question 1 (5 Marks)]
\textbf{Step 1: Augmented Matrix Formulation}
\[
[A|B] = \begin{bmatrix}
1 & 1 & 1 & | & 6 \\
1 & 2 & 3 & | & 10 \\
1 & 2 & \lambda & | & \mu
\end{bmatrix}
\]
\textbf{Step 2: Row Echelon Reduction}
Apply $R_2 \to R_2 - R_1$ and $R_3 \to R_3 - R_1$:
\[
\begin{bmatrix}
1 & 1 & 1 & | & 6 \\
0 & 1 & 2 & | & 4 \\
0 & 1 & \lambda-1 & | & \mu-6
\end{bmatrix}
\]
Apply $R_3 \to R_3 - R_2$:
\[
\begin{bmatrix}
1 & 1 & 1 & | & 6 \\
0 & 1 & 2 & | & 4 \\
0 & 0 & \lambda-3 & | & \mu-10
\end{bmatrix}
\]

\textbf{Step 3: Classification of Solutions}
\begin{enumerate}[label=(\roman*), leftmargin=6mm]
    \item \textbf{Unique Solution:}
    Requires $\text{rank}(A) = \text{rank}[A|B] = 3$ (number of unknowns).
    Hence, the pivot $\lambda - 3 \neq 0 \implies \mathbf{\lambda \neq 3}$, and $\mathbf{\mu \in \mathbb{R}}$ (any value).
    \item \textbf{No Solution (Inconsistent):}
    Requires $\text{rank}(A) = 2$ and $\text{rank}[A|B] = 3$.
    This occurs when $\lambda - 3 = 0$ and $\mu - 10 \neq 0$:
    \[
    \mathbf{\lambda = 3 \quad \text{and} \quad \mu \neq 10}
    \]
    \item \textbf{Infinitely Many Solutions:}
    Requires $\text{rank}(A) = \text{rank}[A|B] = 2 < 3$.
    This occurs when the entire third row becomes zero:
    \[
    \mathbf{\lambda = 3 \quad \text{and} \quad \mu = 10}
    \]
\end{enumerate}

\begin{markingrubric}
\textbf{Mark Distribution:}
$\bullet$ Augmented matrix and row operations: \textbf{2 Marks}\\
$\bullet$ Unique solution condition ($\lambda \neq 3, \mu \in \mathbb{R}$): \textbf{1 Mark}\\
$\bullet$ No solution condition ($\lambda = 3, \mu \neq 10$): \textbf{1 Mark}\\
$\bullet$ Infinite solution condition ($\lambda = 3, \mu = 10$): \textbf{1 Mark}
\end{markingrubric}
\end{solutionbox}

% Solution 2
\begin{solutionbox}[Solution to Question 2 (5 Marks)]
\textbf{Step 1: Auxiliary Equation}
Given $(D^3 + 1)y = 0$, replace $D$ by $m$:
\[
m^3 + 1 = 0
\]
Use algebraic identity $a^3 + b^3 = (a+b)(a^2 - ab + b^2)$:
\[
(m + 1)(m^2 - m + 1) = 0
\]
\textbf{Step 2: Finding Roots}
\begin{itemize}[leftmargin=6mm]
    \item From $m + 1 = 0 \implies \mathbf{m_1 = -1}$.
    \item From $m^2 - m + 1 = 0$, using quadratic formula:
    \[
    m = \frac{-(-1) \pm \sqrt{(-1)^2 - 4(1)(1)}}{2(1)} = \frac{1 \pm \sqrt{-3}}{2} = \mathbf{\frac{1}{2} \pm i\frac{\sqrt{3}}{2}}
    \]
\end{itemize}
\textbf{Step 3: General Solution}
For the real root $m_1 = -1$, the term is $c_1 e^{-x}$.
For complex conjugate roots $\alpha \pm i\beta$ with $\alpha = 1/2, \beta = \sqrt{3}/2$:
\[
e^{\frac{1}{2}x}\left(c_2 \cos\frac{\sqrt{3}}{2}x + c_3 \sin\frac{\sqrt{3}}{2}x\right)
\]
Therefore, the general solution is:
\[
\mathbf{y(x) = c_1 e^{-x} + e^{\frac{x}{2}}\left(c_2 \cos\frac{\sqrt{3}}{2}x + c_3 \sin\frac{\sqrt{3}}{2}x\right)}
\]

\begin{markingrubric}
\textbf{Mark Distribution:}
$\bullet$ Auxiliary equation $m^3 + 1 = 0$ and factorization: \textbf{1.5 Marks}\\
$\bullet$ Correct roots $m = -1, \frac{1}{2} \pm i\frac{\sqrt{3}}{2}$: \textbf{1.5 Marks}\\
$\bullet$ Final general solution with all 3 arbitrary constants: \textbf{2 Marks}
\end{markingrubric}
\end{solutionbox}

% Solution 3
\begin{solutionbox}[Solution to Question 3 (5 Marks)]
\textbf{Step 1: Matrix Inspection \& Row Operations}
\[
A = \begin{bmatrix}
1 & w & r & 0 \\
0 & s & t & 1 \\
-1 & -w & -r & 0 \\
0 & s & t & 1
\end{bmatrix}
\]
Notice that:
\begin{itemize}[leftmargin=6mm]
    \item Row 3 is the exact negative of Row 1: $R_3 = -R_1$.
    \item Row 4 is identical to Row 2: $R_4 = R_2$.
\end{itemize}
\textbf{Step 2: Row Reduction to Echelon Form}
Apply $R_3 \to R_3 + R_1$ and $R_4 \to R_4 - R_2$:
\[
\begin{bmatrix}
1 & w & r & 0 \\
0 & s & t & 1 \\
0 & 0 & 0 & 0 \\
0 & 0 & 0 & 0
\end{bmatrix}
\]
\textbf{Step 3: Determining Rank}
Since $s \neq 0$, the first two rows are non-zero and linearly independent. The remaining two rows are entirely zero.
The number of non-zero rows in the row echelon form is 2.
Therefore:
\[
\mathbf{\text{rank}(A) = 2}
\]

\begin{markingrubric}
\textbf{Mark Distribution:}
$\bullet$ Identifying linear dependency between rows: \textbf{2 Marks}\\
$\bullet$ Performing row operations to echelon form: \textbf{2 Marks}\\
$\bullet$ Correct rank $= 2$: \textbf{1 Mark}
\end{markingrubric}
\end{solutionbox}

% Solution 4
\begin{solutionbox}[Solution to Question 4 (5 Marks)]
\textbf{Step 1: Characteristic Equation of $A$}
For $A = \begin{bmatrix} 1 & 2 \\ 3 & 4 \end{bmatrix}$:
\[
\text{trace}(A) = 1 + 4 = 5, \quad \det(A) = (1)(4) - (2)(3) = 4 - 6 = -2
\]
The characteristic equation is:
\[
\lambda^2 - 5\lambda - 2 = 0
\]
\textbf{Step 2: Verifying Cayley-Hamilton Theorem}
The theorem asserts that $A$ satisfies its own characteristic equation:
\[
A^2 - 5A - 2I = O
\]
Compute $A^2$:
\[
A^2 = \begin{bmatrix} 1 & 2 \\ 3 & 4 \end{bmatrix}\begin{bmatrix} 1 & 2 \\ 3 & 4 \end{bmatrix} = \begin{bmatrix} 1+6 & 2+8 \\ 3+12 & 6+16 \end{bmatrix} = \begin{bmatrix} 7 & 10 \\ 15 & 22 \end{bmatrix}
\]
Now compute $A^2 - 5A - 2I$:
\[
\begin{bmatrix} 7 & 10 \\ 15 & 22 \end{bmatrix} - 5\begin{bmatrix} 1 & 2 \\ 3 & 4 \end{bmatrix} - 2\begin{bmatrix} 1 & 0 \\ 0 & 1 \end{bmatrix} = \begin{bmatrix} 7 - 5 - 2 & 10 - 10 - 0 \\ 15 - 15 - 0 & 22 - 20 - 2 \end{bmatrix} = \begin{bmatrix} 0 & 0 \\ 0 & 0 \end{bmatrix} = O
\]
Hence, Cayley-Hamilton Theorem is verified!

\textbf{Step 3: Calculating $A^{-1}$ using Cayley-Hamilton}
Multiply the verified relation by $A^{-1}$:
\[
A - 5I - 2A^{-1} = O \implies 2A^{-1} = A - 5I \implies A^{-1} = \frac{1}{2}(A - 5I)
\]
\[
A^{-1} = \frac{1}{2}\left(\begin{bmatrix} 1 & 2 \\ 3 & 4 \end{bmatrix} - \begin{bmatrix} 5 & 0 \\ 0 & 5 \end{bmatrix}\right) = \frac{1}{2}\begin{bmatrix} -4 & 2 \\ 3 & -1 \end{bmatrix} = \mathbf{\begin{bmatrix} -2 & 1 \\ 3/2 & -1/2 \end{bmatrix}}
\]

\begin{markingrubric}
\textbf{Mark Distribution:}
$\bullet$ Characteristic equation $\lambda^2 - 5\lambda - 2 = 0$: \textbf{1.5 Marks}\\
$\bullet$ Verification $A^2 - 5A - 2I = O$: \textbf{2 Marks}\\
$\bullet$ Computation of $A^{-1}$: \textbf{1.5 Marks}
\end{markingrubric}
\end{solutionbox}

% Solution 5
\begin{solutionbox}[Solution to Question 5 (5 Marks)]
\textbf{Step 1: Definition and Verification of Hermitian Matrix}
A complex square matrix $A$ is Hermitian if $A^\dagger = (\bar{A})^T = A$.
Given:
\[
A = \begin{bmatrix} 2 & 3+4i \\ 3-4i & 5 \end{bmatrix}
\]
Take the complex conjugate $\bar{A}$:
\[
\bar{A} = \begin{bmatrix} 2 & 3-4i \\ 3+4i & 5 \end{bmatrix}
\]
Transpose $(\bar{A})^T$:
\[
(\bar{A})^T = \begin{bmatrix} 2 & 3+4i \\ 3-4i & 5 \end{bmatrix} = A
\]
Thus, \textbf{$A$ is indeed a Hermitian matrix}.

\textbf{Step 2: Eigenvalues Computation}
Characteristic equation: $\det(A - \lambda I) = 0$:
\[
\begin{vmatrix} 2-\lambda & 3+4i \\ 3-4i & 5-\lambda \end{vmatrix} = 0
\]
\[
(2-\lambda)(5-\lambda) - (3+4i)(3-4i) = 0
\]
\[
(\lambda^2 - 7\lambda + 10) - (3^2 - (4i)^2) = 0
\]
\[
(\lambda^2 - 7\lambda + 10) - (9 + 16) = 0 \implies \lambda^2 - 7\lambda - 15 = 0
\]
Solve for $\lambda$:
\[
\lambda = \frac{7 \pm \sqrt{(-7)^2 - 4(1)(-15)}}{2} = \frac{7 \pm \sqrt{49 + 60}}{2} = \mathbf{\frac{7 \pm \sqrt{109}}{2}}
\]
Since $109 > 0$, both eigenvalues $\mathbf{\lambda_1 = \frac{7 + \sqrt{109}}{2} \approx 8.72}$ and $\mathbf{\lambda_2 = \frac{7 - \sqrt{109}}{2} \approx -1.72}$ are \textbf{strictly real}. This confirms the fundamental theorem that every Hermitian matrix has real eigenvalues.

\begin{markingrubric}
\textbf{Mark Distribution:}
$\bullet$ Proving $(\bar{A})^T = A$ for Hermitian: \textbf{2 Marks}\\
$\bullet$ Characteristic equation $\lambda^2 - 7\lambda - 15 = 0$: \textbf{1.5 Marks}\\
$\bullet$ Finding real eigenvalues: \textbf{1.5 Marks}
\end{markingrubric}
\end{solutionbox}

% Solution 6
\begin{solutionbox}[Solution to Question 6 (5 Marks)]
\textbf{Step 1: Complementary Function ($y_c$)}
Auxiliary equation:
\[
m^2 - 4m + 4 = 0 \implies (m - 2)^2 = 0 \implies m = 2, 2 \text{ (repeated real root)}
\]
\[
\mathbf{y_c = (c_1 + c_2 x) e^{2x}}
\]
\textbf{Step 2: Particular Integral ($y_p$)}
\[
y_p = \frac{1}{(D - 2)^2} e^{2x}
\]
Notice that substituting $D = 2$ gives $(2-2)^2 = 0$ (Case of failure of order 2).
Using the standard operator rule $\frac{1}{(D - a)^n} e^{ax} = \frac{x^n}{n!} e^{ax}$:
Here $a = 2, n = 2$:
\[
y_p = \frac{x^2}{2!} e^{2x} = \mathbf{\frac{1}{2} x^2 e^{2x}}
\]
\textbf{Step 3: General Solution}
\[
y = y_c + y_p = \mathbf{(c_1 + c_2 x) e^{2x} + \frac{1}{2} x^2 e^{2x}}
\]

\begin{markingrubric}
\textbf{Mark Distribution:}
$\bullet$ Complementary function $y_c = (c_1 + c_2 x)e^{2x}$: \textbf{2 Marks}\\
$\bullet$ Case of failure recognition and Particular Integral $y_p$: \textbf{2 Marks}\\
$\bullet$ Final general solution $y = y_c + y_p$: \textbf{1 Mark}
\end{markingrubric}
\end{solutionbox}

\newpage
